\chapter{Reverse Engineering van het Communicatieprotocol}
\label{ch:reverse_engineering}

Dit hoofdstuk vormt de technische kern van het onderzoek. Het beschrijft het proces van het ontsluiten van de gesloten datacommunicatie tussen de 
Laerdal Little Anne QCPR en de Meta Quest 3. Aangezien er geen publieke SDK beschikbaar is voor integratie binnen de Unity-omgeving op Android-basis, 
was een \textit{reverse engineering} traject noodzakelijk om de ruwe sensordata te vertalen naar bruikbare parameters voor de Mixed Reality-omgeving.

\section{Methodologie van het onderzoek}
Het onderzoek naar het protocol werd opgedeeld in drie fasen:
\begin{enumerate}
    \item \textbf{Statische Analyse:} Het decompileren van de officiële Android-applicatie om de broncode te bestuderen.
    \item \textbf{Dynamische Analyse:} Het onderscheppen van live Bluetooth-verkeer tussen de pop en een smartphone via HCI-sniffing.
    \item \textbf{Protocol Mapping:} Het correleren van onderschepte bytes aan fysieke handelingen op de pop en het implementeren van deze logica in C#.
\end{enumerate}

\section{Statische Analyse en de barrière van Obfuscatie}
De eerste poging tot het begrijpen van het protocol bestond uit het decompileren van de officiële Laerdal QCPR-app (APK-bestand) middels tools zoals \textit{Jadx-GUI}. 

Hoewel de decompilatie succesvol was in het herstellen van de mappenstructuur, stuitte het onderzoek direct op \textit{code obfuscation}. 
De fabrikant heeft de broncode behandeld met tools die variabele- en methodenamen vervangen door nietszeggende karakters (bijv. \texttt{a.b(c, d)}). 
Hierdoor was het onmogelijk om de logica achter de Bluetooth-handshake of de data-interpretatie puur uit de broncode af te leiden. Deze doodlopende weg 
bevestigde de noodzaak voor een \textit{black-box} benadering via netwerkanalyse op de hardware-interface.

\begin{figure}[ht]
    \centering
    \includegraphics[width=0.8\textwidth]{jadx_obfuscation.png} % Zorg dat je screenshot zo heet in je graphics map
    \caption{Degecompileerde broncode in Jadx-GUI vertoont zware obfuscatie (variabelen hernoemd naar a, b, c).}
    \label{fig:jadx_obfuscation}
\end{figure}

\section{Dynamische Analyse: HCI Packet Sniffing}
Om de communicatie te begrijpen zonder leesbare broncode, werd besloten om het actieve Bluetooth-verkeer direct te onderscheppen. 

\subsection{Opzet van de testomgeving}
Een standaard smartphone staat het onderscheppen van Bluetooth Low Energy (BLE) pakketten op applicatieniveau meestal niet toe vanwege beveiligingsrestricties. 
Daarom werd een gerouteerd (\textit{rooted}) Android-toestel ingezet. Dit bood de mogelijkheid om de \textit{Bluetooth HCI Snoop Log} te activeren. Deze log legt 
elk pakket vast dat door de Bluetooth-controller van het toestel wordt verwerkt op de \textit{Host Controller Interface} (HCI) laag.

\subsection{Data-acquisitie via Wireshark}
\begin{figure}[ht]
    \centering
    \includegraphics[width=0.6\textwidth]{nrf_connect_uuids.png}
    \caption{Overzicht van de gevonden GATT-services en characteristics in nRF Connect.}
    \label{fig:nrf_uuids}
\end{figure}

Na een reanimatiesessie met de officiële app werd het gegenereerde logbestand (\texttt{btsnoop\_hci.log}) geanalyseerd met \textit{Wireshark}. In deze fase werd 
gezocht naar \textit{GATT (Generic Attribute Profile)} transacties. De focus lag op het identificeren van de \textit{Write Requests} die de initiële verbinding 
tot stand brachten en de \textit{Handle Value Notifications} die de live sensordata bevatten.

\section{Ontcijfering van het Protocol}
Door de Wireshark-logs te correleren met de uiteindelijke Unity-implementatie, zijn de volgende technische specificaties achterhaald:

\subsection{Identificatie van Characteristics}
Er werd vastgesteld dat de pop communiceert via een specifiek Service-UUID-patroon (\texttt{d746-4092-84e7-dad34863fe4a}). De belangrijkste kanalen zijn:
\begin{description}
    \item[Main Service:] \texttt{00000126-d746-4092-84e7-dad34863fe4a}
    \item[Auth Characteristic:] \texttt{000001b1-d746-4092-84e7-dad34863fe4a}
    \item[Wave (Data) Characteristic:] \texttt{00000027-d746-4092-84e7-dad34863fe4a}
\end{description}

\subsection{De "Handshake" en Authentificatie}
Een cruciale ontdekking was dat de pop beveiligd is met een authenticatie-mechanisme. Zonder het versturen van een specifieke \textit{Auth Key} naar de pop,
weigert de hardware sensordata te verzenden. De geanalyseerde handshake bestaat uit drie stappen:
\begin{enumerate}
    \item \textbf{Authenticatie:} Het schrijven van een 20-byte array (startend met \texttt{0x5C 0x0E...}) naar het \textit{Auth-UUID}.
    \item \textbf{Activatie 1:} Een write-opdracht met payload \texttt{0x03} naar UUID \texttt{...012b}.
    \item \textbf{Activatie 2:} Een write-opdracht met payload \texttt{0x03 0xFF 0x00} naar UUID \texttt{...0127}.
\end{enumerate}

\begin{figure}[ht]
    \centering
    \includegraphics[width=1\textwidth]{wireshark_packet.png}
    \caption{Wireshark analyse van een Bluetooth HCI snoop log waarbij de authenticatie-bytes worden verzonden.}
    \label{fig:wireshark_handshake}
\end{figure}

\section{Data Mapping: Van Bytes naar Millimeters}
Na de succesvolle handshake begint de pop met het uitzenden van \textit{notifications} op het \textit{Wave}-kanaal. De data wordt verzonden als een byte-array. 
Uit de implementatie in C\# blijkt dat de relevante informatie over de compressiediepte zich in de laatste byte van het pakket bevindt (\texttt{bytes[bytes.Length - 1]}).

De ruwe waarde in deze byte correleert direct met de fysieke indrukking van de borstkas, waarbij een hexadecimale waarde van \texttt{0x4A} (decimaal 74) 
het maximale meetbereik vertegenwoordigt. In de Unity-omgeving wordt deze waarde genormaliseerd en vertaald naar een procentuele diepte voor de visuele 
feedback-elementen.

\section{Conclusie}
Het succesvol \textit{reverse engineeren} van dit protocol heeft de afhankelijkheid van een officiële SDK weggenomen. Door het omzeilen van de obfuscatie 
via \textit{packet sniffing} en het handmatig implementeren van de handshake in C\#, is een stabiele databrug geslagen tussen de medische hardware en de Meta Quest 3.